\chapter{Running Example}
\label{app:running-example}

Consider the following simple theorem:
\begin{verbatim}
theorem split_imp
(P Q R : Prop)
(hpq : P -> Q)
(hpr : P -> R) :
P -> Q /\ R := by
\end{verbatim}

The initial goal is to prove \(P \to Q \land R\). A local proof action may solve the goal directly:
\begin{verbatim}
intro hp
exact And.intro (hpq hp) (hpr hp)
\end{verbatim}

If Lean accepts this proof body and no goal remains, the current proof state is solved. This action does
not create child proof states. It is a complete local attempt for the current target hole.

A skeleton action may instead expose the two child subgoals \(Q\) and \(R\):
\begin{verbatim}
intro hp
have hq : Q := by
  sorry
have hr : R := by
  sorry
exact And.intro hq hr
\end{verbatim}

This skeleton does more than list two subgoals. It also gives the proof structure that will use the
subgoals after they are solved. The declarations \texttt{hq} and \texttt{hr} are the child subgoals, and
the final line
\begin{verbatim}
exact And.intro hq hr
\end{verbatim}
explains how they are combined to prove the parent goal.

The two placeholders create two child subgoals. When the system focuses on the first child, the scaffold
keeps the surrounding proof structure, but only the proof of \texttt{hq} remains as the target:
\begin{verbatim}
intro hp
have hq : Q := by
  sorry
have hr : R := by
  admit
exact And.intro hq hr
\end{verbatim}

The current node is responsible only for replacing the target \texttt{sorry}. The sibling subgoal
\texttt{hr} is temporarily represented by \texttt{admit}, meaning that it is not solved by this node. In this
scaffold, the proof of \texttt{hq} can be:
\begin{verbatim}
exact hpq hp
\end{verbatim}

When the system focuses on the second child, the roles are reversed:
\begin{verbatim}
intro hp
have hq : Q := by
  admit
have hr : R := by
  sorry
exact And.intro hq hr
\end{verbatim}

Now the target is the proof of \texttt{hr}, and the proof body can be:
\begin{verbatim}
exact hpr hp
\end{verbatim}

The use of \texttt{admit} here is temporary. It allows the system to check one child proof inside the
parent scaffold while the sibling child is still open. These checks are not final proofs of the theorem.
They only verify whether a proposed replacement is valid in the proof context where it will later be used.

After both children are solved, GammaZero inserts the solved proof bodies back into the original skeleton:
\begin{verbatim}
intro hp
have hq : Q := by
  exact hpq hp
have hr : R := by
  exact hpr hp
exact And.intro hq hr
\end{verbatim}

The completed scaffold is then checked again by Lean. This final check is essential: the parent proof is
accepted only if the stitched proof contains no remaining \texttt{sorry} or \texttt{admit} and Lean verifies
the whole block.

This example illustrates the boundary between the two action types. A local proof action tries to close
the current target directly. A skeleton action introduces proof structure and may create
child subgoals. Each child is solved in the scaffold where it was created, and the parent skeleton succeeds
only when all required children are solved and the completed proof is accepted by Lean.

The example is simple, but the same pattern applies to larger proofs. A difficult theorem may require
extra bounds, normalization steps, case-specific facts, or auxiliary lemmas that are not explicit in
the theorem statement. A skeleton makes such hidden structure explicit. Instead of forcing the search to
discover a long tactic sequence step by step, the system can first propose a proof structure, then solve the
child subgoals inside that structure.
